\documentclass[11pt,a4paper]{article}

\usepackage[margin=2.5cm]{geometry}
\usepackage{amsmath}
\usepackage{amssymb}
\usepackage{booktabs}
\usepackage{hyperref}
\usepackage{xcolor}
\usepackage{longtable}
\usepackage{array}
\usepackage{graphicx}
\usepackage{fancyhdr}
\usepackage{titlesec}
\usepackage{multirow}
\usepackage{enumitem}

\definecolor{apexblue}{RGB}{0,51,102}
\definecolor{apexgray}{RGB}{100,100,100}
\definecolor{apexred}{RGB}{180,30,30}

\hypersetup{
  colorlinks=true,
  linkcolor=apexblue,
  citecolor=apexblue,
  urlcolor=apexblue,
  pdftitle={MDD: Credit Valuation Adjustment -- APEX-MDL-0010},
  pdfauthor={Apex Global Bank, Quantitative Research Division}
}

\pagestyle{fancy}
\fancyhf{}
\lhead{\textcolor{apexblue}{\small\textbf{APEX-MDL-0010} \textbar\ Credit Valuation Adjustment (CVA) v1.0}}
\rhead{\textcolor{apexgray}{\small RESTRICTED --- Model Risk Sensitive}}
\lfoot{\textcolor{apexgray}{\small Apex Global Bank \textbar\ Quantitative Research Division}}
\rfoot{\textcolor{apexgray}{\small Page \thepage}}
\renewcommand{\headrulewidth}{0.4pt}
\renewcommand{\footrulewidth}{0.4pt}

\titleformat{\section}{\large\bfseries\color{apexblue}}{{\thesection.}}{0.5em}{}[\titlerule]
\titleformat{\subsection}{\normalsize\bfseries\color{apexblue}}{{\thesubsection}}{0.5em}{}

\setlength{\parskip}{0.5em}
\setlength{\parindent}{0pt}

\begin{document}

\begin{center}
  {\LARGE\bfseries\color{apexblue} Model Development Document}\\[0.4em]
  {\Large\bfseries Credit Valuation Adjustment (CVA)}\\[0.3em]
  {\large Version 1.0}\\[0.6em]
  {\normalsize Apex Global Bank \textbar\ Quantitative Research Division}\\[0.2em]
  {\small\color{apexgray} SR 11-7 Section Reference: Model Development \textbar\ March 2026}
\end{center}

\vspace{1em}
\textcolor{apexred}{\textbf{Classification: RESTRICTED --- Model Risk Sensitive}}
\vspace{1em}
\hrule
\vspace{1em}

% ─────────────────────────────────────────────────────────────────────────────
\section{Document Metadata}
% ─────────────────────────────────────────────────────────────────────────────

\begin{table}[h!]
\centering
\begin{tabular}{@{}ll@{}}
\toprule
\textbf{Field} & \textbf{Value} \\
\midrule
Model Name           & Credit Valuation Adjustment (CVA) \\
Model ID             & APEX-MDL-0010 \\
Version              & 1.0 \\
Status               & In Validation \\
Risk Tier            & Tier 1 \\
Regulatory Framework & Basel III SA-CVA / BA-CVA; IFRS 13 Fair Value; ISDA SIMM \\
Date of Development  & March 2026 \\
Business Owner       & Head of Global Markets Trading (James Okafor) \\
Model Developer      & Quantitative Research --- XVA Team \\
MRO Review           & Pending --- Dr.\ Rebecca Chen \\
CRO Approval         & Pending --- Dr.\ Priya Nair \\
\bottomrule
\end{tabular}
\caption{Document metadata}
\end{table}

\subsection{Version History}

\begin{table}[h!]
\centering
\begin{tabular}{@{}llll@{}}
\toprule
\textbf{Version} & \textbf{Date} & \textbf{Author} & \textbf{Change Summary} \\
\midrule
0.1 & Jan 2026 & XVA Quant Team & Initial draft --- bilateral framework \\
0.9 & Feb 2026 & XVA Quant Team & Added WWR module; stress calibration \\
1.0 & Mar 2026 & XVA Quant Team & Production candidate; submitted for validation \\
\bottomrule
\end{tabular}
\caption{Version history}
\end{table}

% ─────────────────────────────────────────────────────────────────────────────
\section{Executive Summary}
% ─────────────────────────────────────────────────────────────────────────────

\subsection{Model Purpose}

CVA is the market value of counterparty credit risk embedded in a derivatives
portfolio. It represents the expected loss arising from counterparty default,
adjusted for the replacement cost of the derivative at the time of default.
This model computes bilateral CVA (incorporating both counterparty and
own-credit risk) for all OTC derivatives held by Apex Global Bank that are
subject to counterparty credit risk.

CVA is subtracted from the risk-free value of the derivatives portfolio to
arrive at the fair value under IFRS~13. It also feeds the regulatory CVA
capital charge under Basel~III.

\subsection{Business Use Cases}

\begin{enumerate}[noitemsep]
  \item \textbf{Fair Value Accounting (IFRS 13)}: Daily CVA P\&L is reported
        in the income statement. CVA is a component of derivatives fair value
        disclosed in financial statements.
  \item \textbf{Regulatory Capital (Basel III SA-CVA)}: CVA feeds the
        Standardised Approach CVA capital charge, calibrated to market risk
        sensitivities of CVA.
  \item \textbf{Derivatives Pricing}: Front-office traders add CVA when
        pricing new derivatives with uncollateralised or partially
        collateralised counterparties.
  \item \textbf{Counterparty Limit Management}: CVA informs the Credit Risk
        Officer's assessment of exposure relative to credit limits.
  \item \textbf{Hedging}: The CVA desk dynamically hedges CVA sensitivities
        (credit delta CS01, IR delta DV01) using CDS and interest rate
        instruments.
\end{enumerate}

\subsection{Output Description}

\begin{itemize}[noitemsep]
  \item \textbf{Unilateral CVA}: Expected loss from counterparty default only.
  \item \textbf{Bilateral CVA (BCVA)}: Net of CVA and DVA per IFRS~13.
  \item \textbf{CVA Greeks}: Sensitivities to CS01, DV01, FX, and
        volatility --- used for hedging and SA-CVA capital.
  \item \textbf{CVA by counterparty}: Granular attribution for limit
        management.
  \item \textbf{CVA P\&L explain}: Daily change decomposed into market moves,
        credit spread changes, and new trades.
\end{itemize}

\subsection{Key Assumptions}

\begin{enumerate}[noitemsep]
  \item Counterparty default times follow a Poisson process with intensity
        implied by CDS spreads or internal PD estimates.
  \item Recovery rates are fixed at 40\% for senior unsecured counterparties,
        with desk override capability for specific counterparties.
  \item Market risk factors evolve independently of counterparty credit
        quality, except where wrong-way risk is explicitly modelled.
  \item Netting sets are legally enforceable --- ISDA Master Agreement and CSA
        in place and legally reviewed for each counterparty.
  \item Simulation horizon matches the longest trade maturity (up to 30 years
        for long-dated swaps).
\end{enumerate}

\subsection{Known Limitations}

\begin{enumerate}[noitemsep]
  \item Wrong-way risk between FX and sovereign/bank counterparties is
        approximated; full correlation-based WWR simulation is computationally
        expensive.
  \item Recovery rates are deterministic; stochastic recovery adds model
        complexity without materially changing the mean CVA estimate.
  \item The model does not capture re-rating dynamics (gradual credit
        deterioration before default).
  \item Gap risk (jump-to-default over a weekend) is not captured in the
        continuous-time framework.
  \item Proxy CDS spreads for counterparties without observable CDS introduce
        basis risk.
\end{enumerate}

\subsection{Materiality Assessment}

CVA for Apex's derivatives portfolio is currently approximately
\textbf{\$2.1 billion} (gross, pre-netting and collateral). A 1\% change
in average counterparty credit spreads produces approximately
\textbf{\$85 million} in CVA P\&L. This is a Tier~1 model: it feeds both
the income statement and regulatory capital.

% ─────────────────────────────────────────────────────────────────────────────
\section{Business Context and Purpose}
% ─────────────────────────────────────────────────────────────────────────────

\subsection{Business Problem}

Prior to CVA, derivatives were marked to risk-free value, ignoring
counterparty credit quality. A swap with an AAA-rated government counterparty
and an identical swap with a CCC-rated corporate were marked identically ---
a manifest mispricing exposed catastrophically when Lehman Brothers defaulted
with approximately \$35 billion in outstanding OTC derivatives in 2008.

Post-crisis, IFRS~13 (2013) and the Basel~III CVA capital charge (2019)
mandated that banks explicitly value and capitalise counterparty credit risk.

\subsection{Regulatory Requirement}

Under \textbf{Basel~III SA-CVA} (CRR Article 383), banks must hold capital
against the risk that CVA itself moves due to market risk factors. The SA-CVA
approach requires computing CVA sensitivities (delta and vega) to all market
risk factors and aggregating them under a prescribed formula with regulatory
correlations. As Apex does not currently hold IMA-CVA approval, the
Standardised Approach applies to all CVA capital.

\subsection{Model Users}

\begin{table}[h!]
\centering
\begin{tabular}{@{}ll@{}}
\toprule
\textbf{User} & \textbf{Function} \\
\midrule
XVA Desk                   & Daily CVA computation, P\&L, hedging \\
Credit Risk Officer        & Counterparty exposure monitoring \\
Finance / Accounting       & IFRS 13 fair value for financial statements \\
Regulatory Reporting       & Basel III CVA capital charge computation \\
Front Office Pricing Desks & CVA charge in client-facing derivatives prices \\
\bottomrule
\end{tabular}
\caption{Model users}
\end{table}

% ─────────────────────────────────────────────────────────────────────────────
\section{Theoretical Foundation}
% ─────────────────────────────────────────────────────────────────────────────

\subsection{Unilateral CVA}

Under the risk-neutral measure $\mathbb{Q}$, unilateral CVA is:

\begin{equation}
  \text{CVA} = (1 - R) \int_0^T \text{EE}(t)\, d\text{PD}(t)
  \label{eq:cva_continuous}
\end{equation}

where $R$ is the recovery rate, $\text{EE}(t) = \mathbb{E}^{\mathbb{Q}}[\max(V(t), 0)]$
is the Expected Exposure, $V(t)$ is the mark-to-market value of the netting
set at time $t$, and $\text{PD}(t)$ is the cumulative probability of default
by time $t$.

In discrete form over time steps $t_1, t_2, \ldots, t_n$:

\begin{equation}
  \text{CVA} \approx (1 - R) \sum_{i=1}^{n} \text{EE}(t_i)
    \cdot \bigl[\text{PD}(t_{i-1}, t_i)\bigr]
  \label{eq:cva_discrete}
\end{equation}

where $\text{PD}(t_{i-1}, t_i)$ is the marginal probability of default in the
interval $[t_{i-1}, t_i]$.

\subsection{Bilateral CVA}

The bilateral CVA (BCVA) incorporates both counterparty default risk (CVA)
and Apex's own default risk (DVA):

\begin{equation}
  \text{BCVA} = \text{CVA} - \text{DVA}
  \label{eq:bcva}
\end{equation}

The DVA replaces the counterparty's default probability and exposure profile
with Apex's own default probability and the counterparty's exposure to Apex:

\begin{equation}
  \text{DVA} = (1 - R_{\text{Apex}}) \int_0^T \text{ENE}(t)\, d\text{PD}_{\text{Apex}}(t)
  \label{eq:dva}
\end{equation}

where $\text{ENE}(t) = \mathbb{E}^{\mathbb{Q}}[\min(V(t), 0)]$ is the
Expected Negative Exposure.

\subsection{Credit Curve Calibration}

Counterparty default intensities $\lambda(t)$ are bootstrapped from the CDS
spread term structure. Under a flat hazard rate approximation:

\begin{equation}
  s(T) \approx (1 - R) \cdot \lambda
  \label{eq:cds_approx}
\end{equation}

For a flat hazard rate:

\begin{equation}
  \text{PD}(0, T) = 1 - \exp(-\lambda T),
  \qquad \text{where } \lambda = \frac{s}{1 - R}
  \label{eq:flat_hazard}
\end{equation}

For counterparties without observable CDS, proxy mapping assigns the
counterparty to a peer group (sector $\times$ rating $\times$ geography) and
uses the average CDS spread of observable peers. For loan counterparties with
internal ratings, PD is mapped to an implied spread via:

\begin{equation}
  s \approx \frac{\text{PD} \cdot \text{LGD}}{1 - \text{PD} \cdot \text{LGD} \cdot \text{tenor}}
  \label{eq:pd_to_spread}
\end{equation}

\subsection{Exposure Simulation}

Expected Exposure is computed via Monte Carlo simulation:

\begin{enumerate}[noitemsep]
  \item \textbf{Simulate market risk factors} under the risk-neutral measure:
        interest rates (two-factor Hull-White), FX (GBM with local vol), equity
        (GBM with stochastic vol), commodity prices.
  \item \textbf{Re-price all trades} in each netting set on each simulation
        path at each time step.
  \item \textbf{Apply netting}: sum trades within each legally enforceable
        netting set.
  \item \textbf{Apply collateral}: subtract collateral balance accounting for
        MTA, thresholds, and the Margin Period of Risk.
  \item \textbf{Compute EE}: average positive exposures across paths at each
        time step.
\end{enumerate}

Simulation parameters: $N = 10{,}000$ paths (production); $N = 1{,}000$
(intraday); monthly time steps for the first 2 years, quarterly thereafter,
semi-annual beyond 10 years; 2-factor Hull-White interest rate model
calibrated to swaption implied vols; $47 \times 47$ correlation matrix
estimated from 5-year historical data.

\subsection{Wrong-Way Risk}

Wrong-way risk (WWR) arises when counterparty creditworthiness deteriorates
precisely when exposure is highest. For specific WWR, exposure and default
probability are jointly simulated using a bivariate diffusion:

\begin{align}
  d\lambda(t) &= \kappa(\theta - \lambda(t))\, dt + \sigma_\lambda\, dW_\lambda(t)
  \label{eq:wwr_lambda} \\
  dV(t)       &= \mu_V\, dt + \sigma_V\, dW_V(t)
  \label{eq:wwr_v}
\end{align}

with $dW_\lambda \cdot dW_V = \rho_{\text{WWR}}\, dt$, where $\rho_{\text{WWR}}$
is estimated from historical correlation between CDS spreads and derivative
value for the specific counterparty.

\subsection{SA-CVA Capital Calculation}

Under Basel~III SA-CVA, the capital charge is:

\begin{equation}
  K_{\text{SA-CVA}} = m_{\text{CVA}} \sqrt{
    \sum_b \left[\rho \sum_k WS_{k,b}\right]^2
    + (1 - \rho^2) \sum_k WS_{k,b}^2
  }
  \label{eq:sa_cva}
\end{equation}

where $m_{\text{CVA}} = 1.25$ is the supervisory scalar, $\rho$ is the
prescribed intra-bucket correlation, and $WS_{k,b}$ is the weighted
sensitivity of CVA to risk factor $k$ in bucket $b$. Sensitivities are
computed by bumping each risk factor $\pm 1$bp and recomputing CVA.

\subsection{Literature References}

\begin{itemize}[noitemsep]
  \item Gregory, J.\ (2012). \textit{Counterparty Credit Risk and Credit Value
        Adjustment}. Wiley Finance.
  \item Pykhtin, M., Zhu, S.\ (2007). ``A Guide to Modelling Counterparty
        Credit Risk.'' \textit{GARP Risk Review}, Jul/Aug 2007.
  \item BCBS (2015). \textit{Review of the Credit Valuation Adjustment Risk
        Framework}. Basel Committee on Banking Supervision.
  \item Brigo, D., Capponi, A.\ (2010). ``Bilateral Counterparty Risk with
        Application to CDSs.'' \textit{Risk Magazine}, March 2010.
\end{itemize}

% ─────────────────────────────────────────────────────────────────────────────
\section{Data Requirements and Governance}
% ─────────────────────────────────────────────────────────────────────────────

\subsection{Input Data Sources}

\begin{longtable}{@{}llll@{}}
\toprule
\textbf{Data Item} & \textbf{Source} & \textbf{Frequency} & \textbf{Fallback} \\
\midrule
\endhead
\bottomrule
\endfoot
CDS spread term structures   & Bloomberg CDSW / Markit & Daily EOD   & Prior day + spread adjustment \\
Interest rate curves (SOFR, ESTR) & Bloomberg / Internal & Intraday & Prior EOD curve \\
FX spot and forward rates    & WM/Reuters 4pm fix       & Daily       & Bloomberg BGN \\
Equity prices                & Exchange feeds           & Real-time   & Prior close \\
Implied volatility surfaces  & Bloomberg OVDV           & Daily       & Prior day + ATM adj. \\
Trade / netting set data     & Front office (Murex)     & Real-time   & Last confirmed booking \\
CSA / collateral terms       & ISDA MA database         & On amendment & Legal review required \\
Recovery rate assumptions    & Credit risk database     & Quarterly   & 40\% senior unsecured \\
Historical correlations      & Risk data warehouse      & Monthly     & 3-year lookback minimum \\
\end{longtable}

\subsection{Data Quality Requirements}

\begin{itemize}[noitemsep]
  \item CDS completeness: $\geq 95\%$ of counterparty PDs from observable CDS
        or proxy; $\leq 5\%$ from internal PD mapping.
  \item Proxy staleness: peer group average CDS age $>5$ business days triggers
        alert.
  \item Trade completeness: 100\% of trades loaded to CVA system within T+1
        of execution; same-day for trades $>\$100$M notional.
  \item CSA terms: any amendment reflected in CVA system within 2 business days
        of legal execution.
\end{itemize}

\subsection{BCBS 239 Data Lineage}

Every CVA output must be traceable to: (1) underlying trade data, (2) market
data snapshot (time-stamped), (3) model version tag, and (4) netting set
assignment with legal review date. Lineage is captured in the risk data
warehouse with immutable audit log. Full lineage for any CVA figure is
reproducible within 4 hours of a regulatory request.

% ─────────────────────────────────────────────────────────────────────────────
\section{Methodology and Implementation}
% ─────────────────────────────────────────────────────────────────────────────

\subsection{Production Run Schedule}

\begin{table}[h!]
\centering
\begin{tabular}{@{}llll@{}}
\toprule
\textbf{Run Type} & \textbf{Timing} & \textbf{Paths} & \textbf{Purpose} \\
\midrule
Intraday indicative & Every 2 hours & 1,000  & Trader pricing / limit monitoring \\
EOD official        & 18:00 EST     & 10,000 & IFRS 13 P\&L, regulatory reporting \\
Weekly stressed     & Saturday      & 20,000 & Stressed CVA for ICAAP \\
Scenario (on demand) & As needed    & 5,000  & New trade pre-approval \\
\bottomrule
\end{tabular}
\caption{Production run schedule}
\end{table}

\subsection{Computational Architecture}

The EOD run involves approximately $4.5 \times 10^9$ pricing evaluations
($10{,}000$ paths $\times$ 30 time steps $\times$ $\sim$15,000 active trades).
Implementation uses GPU acceleration (NVIDIA A100), vectorised
\texttt{numpy}/\texttt{cupy} for exposure aggregation, parallelisation by
netting set, and first-order Taylor approximation for trades below \$5M
notional. EOD runtime target: $\leq 45$ minutes.

\subsection{Netting and Collateral Logic}

\begin{equation}
  \text{Exposure}(t, \omega) = \max\!\left(\sum_i V_i(t, \omega) - C(t, \omega),\; 0\right)
  \label{eq:exposure}
\end{equation}

where the collateral balance evolves as:

\begin{equation}
  C(t) = C(0) + \sum_{\text{margin calls settled up to }(t - \tau)} \Delta C
  \label{eq:collateral}
\end{equation}

with $\tau = 10$ business days (MPoR for daily-margined counterparties;
20 days for weekly-margined). Minimum Transfer Amount (MTA) and thresholds
are applied per CSA terms.

\subsection{CVA Hedging}

The CVA desk actively hedges:
\begin{itemize}[noitemsep]
  \item \textbf{CS01}: via single-name CDS or CDS indices (CDX IG, iTraxx
        Main);
  \item \textbf{DV01}: via interest rate swaps;
  \item \textbf{FX delta}: via FX forwards.
\end{itemize}

Residual unhedged CVA (illiquid counterparties without liquid CDS) is
disclosed separately in risk reports.

% ─────────────────────────────────────────────────────────────────────────────
\section{Model Testing and Backtesting}
% ─────────────────────────────────────────────────────────────────────────────

\subsection{Proxy Backtesting (P\&L Explain)}

CVA P\&L is compared against Greeks-predicted P\&L. Unexplained CVA P\&L
exceeding 10\% of total for more than 5 days in a quarter triggers a model
review.

\subsection{EE Profile Validation}

The model-generated EE profile is compared to:
\begin{enumerate}[noitemsep]
  \item Independent analytical solution (closed-form for simple vanilla IRS);
  \item ORE (Open Source Risk Engine, QuantLib-based) benchmark.
\end{enumerate}

Agreement thresholds: EE profiles $\leq 3\%$ mean absolute deviation; CVA
point estimate $\leq 5\%$ relative deviation; CS01 $\leq 5\%$ relative
deviation.

\subsection{Stress Testing}

\begin{table}[h!]
\centering
\begin{tabular}{@{}lrrr@{}}
\toprule
\textbf{Scenario} & \textbf{CDS Shock} & \textbf{IR Shock} & \textbf{CVA Impact} \\
\midrule
IG counterparties widen 100bp  & +100bp IG  & 0      & +\$420M \\
HY counterparties widen 500bp  & +500bp HY  & 0      & +\$180M \\
Correlation breakdown ($\rho \to 1$) & 0    & 0      & +\$340M \\
WWR activation (top 10 cptys)  & +200bp     & 0      & +\$290M \\
Combined GFC analog            & +300/+800bp & +200bp & +\$1.1B \\
\bottomrule
\end{tabular}
\caption{Stress scenario CVA impacts (estimated, current portfolio)}
\end{table}

% ─────────────────────────────────────────────────────────────────────────────
\section{Performance Metrics and Calibration Standards}
% ─────────────────────────────────────────────────────────────────────────────

\begin{longtable}{@{}llll@{}}
\toprule
\textbf{Metric} & \textbf{Acceptable Range} & \textbf{Frequency} & \textbf{Action if Breached} \\
\midrule
\endhead
\bottomrule
\endfoot
CVA P\&L Explain         & $>85\%$           & Daily     & Investigate within 2 business days \\
EE Benchmark Deviation   & $<5\%$ relative   & Quarterly & Formal model review within 30 days \\
CS01 Accuracy            & $<5\%$ relative   & Monthly   & Desk notification; review within 15 days \\
EOD Computation Time     & $<45$ minutes     & Daily     & Technology escalation \\
CDS Data Completeness    & $>95\%$           & Daily     & Data quality alert; proxy review \\
\end{longtable}

\subsection{Recalibration Schedule}

\begin{itemize}[noitemsep]
  \item \textbf{Correlation matrix}: Monthly, rolling 5-year window.
  \item \textbf{Credit curves}: Daily (automated from Bloomberg/Markit feeds).
  \item \textbf{Volatility surfaces}: Daily (automated).
  \item \textbf{Recovery rate assumptions}: Quarterly review by Credit Risk Officer.
  \item \textbf{WWR parameters}: Semi-annual, or immediately after a
        significant counterparty credit event.
\end{itemize}

\subsection{Failure Criteria}

The model is considered to have failed if any of the following occur:
\begin{enumerate}[noitemsep]
  \item CVA P\&L explain $< 70\%$ for three consecutive business days.
  \item EE benchmark deviation $> 15\%$ on a full portfolio comparison.
  \item EOD run fails to complete within 2 hours (SLA breach).
  \item Data completeness drops below 90\% for more than 2 consecutive days.
  \item Model produces negative CVA for a counterparty without contractual
        explanation.
\end{enumerate}

% ─────────────────────────────────────────────────────────────────────────────
\section{Limitations, Assumptions, and Known Weaknesses}
% ─────────────────────────────────────────────────────────────────────────────

\subsection{Fixed Recovery Rate}

Recovery is fixed at 40\% for senior unsecured counterparties. Recovery rates
are highly variable in practice (20--70\% depending on industry and economic
cycle; approximately 28\% during 2008--2009).

\textbf{Compensating control}: Recovery floored at 35\% in stressed scenario
runs. Desk override for specific counterparties.

\subsection{Independent Defaults (No Contagion)}

The model treats counterparty defaults as driven by idiosyncratic factors only.
In systemic events, financial institution counterparties become highly
correlated.

\textbf{Compensating control}: ``Correlation Breakdown'' stress scenario
($\rho \to 1$) bounds this risk. Counterparty concentration limits further
reduce tail exposure.

\subsection{Static Netting Set}

The netting set composition is fixed at the valuation date. New trades added
after valuation are not reflected until the next run.

\textbf{Compensating control}: Same-day booking requirement for trades
$> \$100$M notional, with manual CVA adjustment same day.

\subsection{Linear Exposure Approximation for Small Trades}

Trades below \$5M notional use a first-order Taylor approximation to EE
(approximately 40\% of trades by count; 2\% by notional). Analysis shows
approximation error $< 1\%$ of CVA contribution for this segment.

\subsection{Proxy CDS Spread Risk}

Approximately 8\% of counterparties (3\% by CVA contribution) lack observable
CDS. Proxy basis risk is unobservable by construction.

\textbf{Compensating control}: 25\% add-on to CVA for all proxy-mapped
counterparties. Proxy mapping reviewed quarterly.

% ─────────────────────────────────────────────────────────────────────────────
\section{Compensating Controls and Risk Mitigants}
% ─────────────────────────────────────────────────────────────────────────────

\begin{longtable}{@{}lll@{}}
\toprule
\textbf{Risk} & \textbf{Control} & \textbf{Owner} \\
\midrule
\endhead
\bottomrule
\endfoot
WWR underestimation       & WWR module; bivariate diffusion; stress scenarios & XVA Quant \\
Recovery rate error       & Conservative 35\% floor; quarterly review         & Credit Risk \\
Proxy CDS spread error    & 25\% CVA add-on for proxy counterparties           & XVA Quant \\
Correlation instability   & Correlation breakdown stress scenario              & Market Risk \\
Booking lag               & Same-day booking for trades $>\$100$M             & Operations \\
Computation failure       & Prior-day CVA + Greeks $\times$ moves fallback    & XVA Technology \\
Contagion / cascade risk  & Concentration limits; stress testing               & Credit Risk \\
\end{longtable}

\subsection{Model Reserve}

A model reserve of \textbf{\$45 million} is held against CVA model uncertainty:
\begin{itemize}[noitemsep]
  \item Proxy spread basis risk: \$15M.
  \item Wrong-way risk model simplification: \$20M.
  \item Recovery rate uncertainty: \$10M.
\end{itemize}

% ─────────────────────────────────────────────────────────────────────────────
\section{Change Management}
% ─────────────────────────────────────────────────────────────────────────────

\textbf{Material changes} requiring full MRO re-validation: change in
simulation methodology (number of risk factors, correlation structure), change
in collateral treatment or MPoR assumption, change in WWR model specification,
recalibration that changes portfolio CVA by more than 5\%, or addition of new
product types to CVA scope.

\textbf{Non-material changes} via expedited review (MRO notification within
5 business days): new counterparty added, CDS proxy group reassignment,
computational optimisation with no methodology change.

\textbf{Emergency patch procedure}: interim CVA estimate via prior-day CVA +
Greeks $\times$ moves; root cause identified within 4 hours; patch deployed
and validated by XVA team lead; MRO notified same day; full incident report
within 5 business days.

% ─────────────────────────────────────────────────────────────────────────────
\section{Approvals and Sign-Off}
% ─────────────────────────────────────────────────────────────────────────────

\begin{table}[h!]
\centering
\begin{tabular}{@{}llll@{}}
\toprule
\textbf{Role} & \textbf{Name} & \textbf{Status} & \textbf{Date} \\
\midrule
Model Developer       & XVA Quant Team      & $\checkmark$ Submitted    & March 2026 \\
Model Risk Officer    & Dr.\ Rebecca Chen   & $\circ$ Pending Validation & --- \\
Chief Risk Officer    & Dr.\ Priya Nair     & $\circ$ Pending MRO Approval & --- \\
Business Owner        & James Okafor        & $\circ$ Pending CRO Approval & --- \\
Regulatory Capital    & Finance / Reg.\ Reporting & $\circ$ Pending    & --- \\
\bottomrule
\end{tabular}
\caption{Approval status}
\end{table}

\vspace{2em}
\hrule
\vspace{0.5em}
{\small\color{apexgray}
Document Classification: RESTRICTED --- Model Risk Sensitive\\
Apex Global Bank \textbar\ Quantitative Research Division \textbar\ XVA Team\\
Next scheduled review: March 2027 (or upon material model change)
}

\end{document}
